\documentclass[12pt]{article}
\usepackage[utf8]{inputenc}
\usepackage[T1]{fontenc}
\usepackage{hyperref}
\title{Kaqchikel, Northern Language Resources for K-12 Curriculum Development}
\author{Compiled by Claude Sonnet 5.5 (claude-sonnet-5-5) for Andr\'as Kornai}
\date{}
\begin{document}
\maketitle

\section{Introduction}
Northern Kaqchikel (ISO 639-3 ckc) is a Mayan language of Guatemala. The search results covered Kaqchikel in general and did not distinguish the Northern variety; whether the items below describe ckc specifically is not established. The identification line in the request gives the family as ``Bookkeeping'', which is a Glottolog bookkeeping label, not a language family. The searches were limited and the report is short.

\section{Primary Educational Resources}
\begin{itemize}
\item \emph{La \"utz aw\"ach? Introduction to Kaqchikel Maya Language} (University of Texas Press, 2006), beginner textbook with exercises, grammar summary and glossary: \url{https://utpress.utexas.edu/9780292714601/la-utz-awach/}
\item \emph{Kik'aslemal ri Kaqchikela'} (Little, Maxwell, Kawoq), advanced textbook with Kaqchikel/Spanish/English glossary (retailer listing): \url{https://northtownbooks.com/book/9781477333426}
\end{itemize}

\section{Linguistic and Reference Materials}
No standalone comprehensive dictionary was found in the results. The textbooks above include glossaries and grammar appendices. The Native Languages site lists further Kaqchikel materials: \url{https://www.native-languages.org/kaqchikel.htm}

\section{Religious and Community Resources}
Nothing relevant was found in the search results.

\section{Digital and Audio Resources}
The Native Languages page above mentions an online Spanish-Kaqchikel glossary, some audio clips and a Kaqchikel Mayan Resource Center; those targets were not individually seen as URLs and are not cited.

\section{Curriculum Development Context}
No school curriculum, ministry or bilingual-education documents were found. The available textbooks are aimed at adult second-language learners. Developers should check Guatemalan bilingual education sources directly.

\end{document}